\scanpage{389}\mspage{128}

tenir en \unsure{se ramenant|se restreignant|raisonnant} tout de suite, pour la théorie des opérations $\Omega$, aux complexes de Kan, \deleted{\unsure{Mais}} et aux chemins de type $\Delta_+$. Mais pour les composer \uline{fonctoriellement}, il me semble y avoir toujours le même problème -- on trouve des chemins de type $\Delta_+\Delta_+$, \inserted{$=\Delta_2^{\unsure{+|\to}}$} et non pas $\Delta_+$, -- mais on peut alors \deleted{\unsure{contracter, pour obtenir}} \inserted{composer les homotopies} directes (revient à s'induire sur $\Delta_+\simeq\{0,2\}\subset\Delta_2^+$).

Or j'ai fait le \unsure{formalisme} des catégories de chemins sans me soucier des conditions \unsure{du|de} type Kan. Je présume que dans le cadre semi-simplicial, les $\underline{\mathrm{Hom}}(\Delta_1^{\unsure{+|\to}},X)$ ne donnent \uline{pas} le \unsure{bon} objet de chemins, pour $X$ pas de Kan. (Car sinon, il semblerait que pour $X=X_{\unsure{\bullet|0}}$ à composantes \uline{finies}, le calcul des $\pi_{\unsure{i|1}}$ devrait être fini, ce qui n'est pourtant pas le cas, m'a-t-on toujours assuré \ldots)

Je présume pourtant, quand on prend $\overline\Sigma=\mathrm{Fl}(X)$, que les \deleted{catégories} objets $\underline{\mathrm{Ch}}_\infty(X;a,b)$ et $\underline{\mathrm{Ch}}_\infty(\overline X;a,b)$ \deleted{\unsure{-}} sont \deleted{\unsure{pas}} \inserted{\deleted{en général}} isomorphes dans $\mathrm{Hot}$, quand on immerge $X$ dans une enveloppe de Thomason\unsure{-| }Kan $\overline X$. \deleted{Cela \unsure{me mènerait à|me ferait} \unsure{dire|donc} que la question d'associativité pourtant,}

% archive sheet number: 388
